\documentclass[10pt,twocolumn,letterpaper]{article}

\usepackage[review,pagenumbers]{cvpr}
%% This file contains a number of tweaks that are typically applied to the main document.
%% They are not enabled by default, but can be enabled by uncommenting the relevant lines.

%%
%% Inline annotations; for predefined colors, refer to "dvipsnames" in the xcolor package:
%% https://tinyurl.com/overleaf-colors
%%
\newcommand{\red}[1]{{\color{red}#1}}
\newcommand{\todo}[1]{{\color{red}#1}}
\newcommand{\TODO}[1]{\textbf{\color{red}[TODO: #1]}}
%%
%% disable for camera ready / submission by uncommenting these lines  
%%
% \renewcommand{\TODO}[1]{}
% \renewcommand{\todo}[1]{#1}

%%
%% work harder in optimizing text layout. Typically shrinks text by 1/6 of page, enable
%% it at the very end of the writing process, when you are just above the page limit
%%
% \usepackage{microtype}

%%
%% fine-tune paragraph spacing
%%
% \renewcommand{\paragraph}[1]{\vspace{.5em}\noindent\textbf{#1.}}

%%
%% globally adjusts space between figure and caption
%%
% \setlength{\abovecaptionskip}{.5em}


%%
%% Allows "the use of \paper to refer to the project name"
%% with automatic management of space at the end of the word
%%
% \usepackage{xspace}
% \newcommand{\paper}{ProjectName\xspace}

%%
%% Commonly used math definitions
%%
% \DeclareMathOperator*{\argmin}{arg\,min}
% \DeclareMathOperator*{\argmax}{arg\,max}

%%
%% Tigthen underline
%%
% \usepackage{soul}
% \setuldepth{foobar}

\definecolor{cvprblue}{rgb}{0.21,0.49,0.74}
\usepackage[pagebackref,breaklinks,colorlinks,allcolors=cvprblue]{hyperref}

\def\paperID{*****}
\def\confName{CVPR}
\def\confYear{2026}

\title{Monocular Metric Depth with Video Consistency}

\author{Yuming Zhang\\
AIAA 3201, Spring 2027\\
50025933
\and
Jiayang Liu\\
AIAA 3201, Spring 2027\\
50027216
}

\begin{document}
\maketitle

\begin{abstract}
We estimate dense depth from one RGB image or a short monocular video, put that depth in meters, and keep it stable over time.
The image benchmarks are Middlebury Stereo~v3 and the official NYU Depth~v2 test split.
StereoSGBM converts disparity to metric depth with the supplied calibration.
Depth Anything~V2-Small is a relative baseline and is scale-and-shift aligned only for scoring.
UniDepthV2-Small predicts metric depth from the published intrinsics, with no test-time alignment.
A known standing height then rescales a whole image, and the same depth ranks indoor objects and raises a 1.5\,m proximity warning.
Flow-guided fusion is run on two self-recorded indoor clips.
Code: \url{https://github.com/Oliver-ming/Intro_to_CV_project_3}.
\end{abstract}

\section{Introduction}
A single photograph does not determine metric depth.
The same image is explained by a small nearby scene or a large distant one, so a network that is only trained to rank pixels can look right after a test-time fit and still be useless for measuring distance.
Video adds a second failure: the fit can change from frame to frame, which shows up as flicker when the camera is still and as scale drift when it moves.

This project keeps those two ideas separate.
Relative depth may be aligned to a benchmark.
Metric depth may not.
The system we evaluate for applications is UniDepthV2 after an optional known-size scale, not the aligned relative map.
The course asks for three parts.
Part~1 is classical stereo and a frame-wise relative model.
Part~2 is a zero-shot metric model.
Part~3 is a lightweight change with an ablation: known-size calibration, plus a depth application, with flow fusion ready for the self-recorded clips.

\section{Related Work}
Semi-global matching turns a rectified pair into a disparity image by optimizing a path-wise energy~\cite{hirschmuller2008sgm}.
KITTI made outdoor stereo and laser depth a standard driving benchmark~\cite{geiger2012kitti}.
We follow the course datasets instead: Middlebury's quarter-resolution training pairs for stereo~\cite{scharstein2014middlebury}, and the labeled NYU Depth~v2 split for indoor monocular depth~\cite{silberman2012nyu}.

Dense prediction transformers replaced fully convolutional depth heads~\cite{ranftl2021dpt}.
Depth Anything trains that family on large-scale unlabeled images~\cite{yang2024depthanything}, and Depth Anything~V2 improves the teacher and the synthetic supervision~\cite{yang2024dav2}.
The V2-Small checkpoint we use is relative: its output is defined up to a scale and a shift.
UniDepth and UniDepthV2 predict metric depth and camera intrinsics together, so a published camera can be passed in at inference~\cite{piccinelli2024unidepth,piccinelli2025unidepthv2}.
Depth Pro targets sharp metric depth in under a second and defines the scale-invariant boundary F1 we report~\cite{bochkovskii2025depthpro}.
RAFT is the optical-flow model used to warp depth between frames~\cite{teed2020raft}.
Video Depth Anything keeps a relative depth consistent over long videos~\cite{chen2025vda}; that comparison is blocked on the self-recorded clips.

Known object height is an old cue for metric scale.
ScaleNet uses an adult height of about 1.70\,m and a car height of about 1.59\,m, and it keeps a person only when the head and both ankles are visible~\cite{zhu2020scalenet}.
We use those priors, and we treat a one-standard-deviation shift of them as the noisy-cue ablation.

\section{Method}
\subsection{Stereo baseline}
On each Middlebury training pair we run StereoSGBM and a left-right consistency check.
Invalid pixels are those whose forward and reverse disparities disagree by more than one pixel, or that have no match.
Metric depth uses the calibration file,
\begin{equation}
Z = \frac{f B}{d + \mathrm{doffs}},
\end{equation}
with the baseline converted from millimeters to meters.
A one-pixel disparity error then moves depth by $Z^2/(fB)$.

\subsection{Frame-wise relative depth}
Depth Anything~V2-Small runs independently on each NYU test image.
For the table only, we solve a per-image scale and shift in depth space by least squares on the Eigen crop.
A median scale with no shift is reported as a negative control.
Neither number is called metric depth.

\subsection{Metric depth}
UniDepthV2-Small receives the NYU Kinect intrinsics at inference and returns meters.
We do not fit the test set.
The same camera is the one used to turn a box height into meters,
\begin{equation}
h = \frac{p_y \, Z}{f_y},
\end{equation}
where $p_y$ is the vertical side of the box and $Z$ is the median predicted depth inside it.

\subsection{Known-size scale}
A person or a car may vote only if the detector score is at least 0.5, the box does not touch the image border, and, for a person, the nose and both ankles are visible.
The size prior is 1.70\,m for an adult and 1.59\,m for a car.
The image scale is the median of prior over implied height.
An image with no vote keeps the network scale.
Because a single NYU image is a one-frame clip, that median is the clip scale.
The noisy ablation repeats the vote at $+1\sigma$ and $-1\sigma$ ($0.09$\,m and $0.21$\,m).

\subsection{Application}
Indoor obstacles are detector boxes for a person, chair, couch, bed, dining table, toilet, television, refrigerator, or potted plant.
The distance of a box is the median metric depth of the pixels inside it.
Ranking is the fraction of box pairs whose nearer-or-farther order matches the ground-truth medians, plus the absolute relative error of those medians.
A frame warns if the nearest obstacle is closer than 1.5\,m.
Precision, recall, and F1 compare that decision with the same rule on ground truth.
The point cloud lifts the same depth with the same intrinsics. It is a view, not a score.

\subsection{Video, not yet run}
RAFT-Small warps the previous fused depth onto the current frame.
Pixels whose forward and backward flow disagree keep the current UniDepth prediction.
The other pixels mix the two maps with a weight that falls as the flow residual grows.
The first frame is UniDepth alone.
Temporal error is the mean relative gap after this warp, on pixels with positive depth.
Static-camera and moving-camera clips will be reported separately.
Those clips are not in the repository yet.

\section{Experiments}
\subsection{Protocol}
Middlebury uses the 15 quarter-resolution training pairs and the non-occluded mask.
NYU uses the official 654-image test indices in \texttt{splits.mat}, the Eigen crop $[45{:}471,\,41{:}601]$, and depths from $10^{-3}$\,m to $10$\,m.
Depth metrics follow the Depth Anything~V2 \texttt{eval\_depth} definitions: AbsRel, SqRel, RMSE, RMSE-log, and $\delta<1.25^{k}$.
Boundary quality is Depth Pro's scale-invariant boundary F1 on that same crop.
We do not score KITTI or a sparse laser, and we do not score an RGB edge.

\subsection{Stereo}
Table~\ref{tab:stereo} summarizes StereoSGBM.
The mean fraction of non-occluded pixels with disparity error above 1 pixel is 0.368.
Textured, high-disparity scenes (Jadeplant, PianoL, Vintage) are worse.
On the Adirondack calibration ($f=1038$\,px, $B=0.176$\,m), a one-pixel disparity error is 5.5\,mm at 1\,m and 350\,mm at 8\,m.
Disparity noise is therefore a near-field measurement and a far-field guess.

\begin{table}[t]
\centering
\small
\begin{tabular}{lcc}
\hline
Scene group & bad-1.0 & left-right kept \\
\hline
Best three (Teddy, Recycle, Adirondack) & 0.25--0.26 & 0.73 \\
Mean of 15 training pairs & 0.368 & -- \\
Jadeplant / PianoL / Vintage & 0.50--0.58 & 0.38--0.48 \\
\hline
\end{tabular}
\caption{StereoSGBM on Middlebury 2014 quarter-resolution training pairs. bad-1.0 is the fraction of non-occluded pixels with disparity error above 1\,px.}
\label{tab:stereo}
\end{table}

\subsection{NYU depth}
Table~\ref{tab:nyu} compares the relative and metric frame-wise models.
Scale-and-shift alignment is required for V2-Small.
Median scaling alone raises AbsRel from 0.132 to 1.35, which is the unknown-shift failure the method section separates from metric depth.
UniDepthV2-Small, with the dataset camera and no alignment, reaches AbsRel 0.093 and $\delta_1$ 0.925.

\begin{table}[t]
\centering
\small
\begin{tabular}{lccc}
\hline
Method & AbsRel & RMSE & $\delta_1$ \\
\hline
V2-Small, scale and shift & 0.132 & 0.427 & 0.845 \\
V2-Small, median scale only & 1.348 & 3.492 & 0.191 \\
UniDepthV2-Small, no alignment & 0.093 & 0.380 & 0.925 \\
\hline
\end{tabular}
\caption{Official NYU Depth~v2 test split, Eigen crop, 10\,m cap. The first two rows are aligned relative depth. The third row is metric.}
\label{tab:nyu}
\end{table}

\subsection{Boundary, scale, and the application}
Table~\ref{tab:part3} is the Part~3 result on the same 654 images.
A global scale does not change inverse-depth ratios, so the boundary F1 of UniDepth stays 0.117 after calibration.
The aligned relative model is slightly higher, at 0.124.
Only 60 of 654 images cast a scale vote.
On the full test set the clean prior raises AbsRel from 0.093 to 0.273, and both one-sigma shifts stay near 0.26--0.29.
The prior is the wrong tool on a model that is already metric: a detector box mixes background into the median, a sitting or partly occluded person still trips the ankle test, and one bad ratio rescales the whole image.
The course asks for this analysis when the third part does not improve the main metric.

The application still reads the calibrated depth, because that is the system depth defined above.
Across 2{,}582 obstacle boxes and 9{,}127 pairs, order accuracy is 0.932 and obstacle AbsRel is 0.324.
The 1.5\,m warning, on 163 frames whose nearest ground-truth obstacle is inside the threshold, has precision 0.911, recall 0.693, and F1 0.788.
Order is easier than meters, which matches the damage done by the global scale.
UniDepthV2-Small itself takes 46\,ms per $480\times640$ frame in this loop.
The process that also holds the detector and V2-Small peaks at about 1.0\,GiB.

\begin{table}[t]
\centering
\small
\begin{tabular}{lc}
\hline
Quantity & Value \\
\hline
UniDepth boundary F1 & 0.117 \\
Calibrated boundary F1 & 0.117 \\
V2-Small aligned boundary F1 & 0.124 \\
Calibrated AbsRel & 0.273 \\
$+1\sigma$ prior AbsRel & 0.290 \\
$-1\sigma$ prior AbsRel & 0.256 \\
Images with a scale vote & 60 / 654 \\
Pairwise order accuracy & 0.932 \\
Obstacle AbsRel & 0.324 \\
Proximity P / R / F1 & 0.911 / 0.693 / 0.788 \\
\hline
\end{tabular}
\caption{NYU test, Part~3. Boundary F1 is scale-invariant. The proximity threshold is 1.5\,m. AbsRel rows use the Eigen crop. Application rows use detector boxes.}
\label{tab:part3}
\end{table}

\subsection{Self-recorded clips}
Two indoor recordings were transcoded to $1280\times720$ at 15\,fps.
\texttt{still.mp4} keeps the camera fixed for 19.9\,s while a person walks.
\texttt{moving.mp4} moves the camera for 19.4\,s while the person stands.
Both show the head and both feet, plus a chair in the near field.
A 4.8\,s moving clip was not used.
UniDepthV2-Small with RAFT fusion gives a flow-warped temporal AbsRel of 0.021 on the still clip and 0.041 on the moving clip.
The colored depth video uses one range for the whole clip.
There is no ground truth, so these clips have no drift score and no metric AbsRel.
Video Depth Anything is not in this comparison yet.

\section{Conclusion}
Stereo makes the $Z^2$ growth of disparity error concrete.
On NYU, a relative model needs a scale and a shift before it matches ground truth, and a metric model does not.
Known-size calibration does not improve NYU AbsRel, and the boundary score is unchanged by a global scale.
The open measurements are the temporal numbers on the two phone clips.
The clips are also where a person entering mid-video would have changed a per-frame scale; the implementation uses one scale for the whole clip so that entry does not cut the depth.

{
\small
\bibliographystyle{ieeenat_fullname}
\bibliography{main}
}

\end{document}
